\documentclass[10pt,a4paper]{amsart}
\usepackage[margin=19mm]{geometry}
\usepackage{amsmath,amssymb,mathtools}
\usepackage[scr=rsfs,cal=euler]{mathalfa}
\usepackage{xfrac}
\usepackage[unicode,hidelinks,bookmarks=false]{hyperref}
\newcommand{\ie}{i.e.\ }
\newcommand{\cf}{cf.\ }
\newcommand{\resp}{respectively\ }
\newcommand{\Subsection}{\subsection}
\providecommand{\nobreakdash}{\nobreak\hbox{-}}
\setlength{\emergencystretch}{2em}
\multlinegap0pt
\usepackage{bm}
\usepackage{bbm}
\usepackage{dsfont}
\usepackage{xspace}
\usepackage{cancel}
\usepackage{mathtools}
\usepackage{amsthm}
\makeatletter
\@ifundefined{thm}{\newtheorem{thm}{Thm}}{}
\@ifundefined{lemma}{\newtheorem{lemma}[thm]{Lemma}}{}
\@ifundefined{prop}{\newtheorem{prop}[thm]{Proposition}}{}
\makeatother
\newcommand{\bA}{\mathbf{A}}
\newcommand{\bB}{\mathbf{B}}
\newcommand{\bD}{\mathbf{D}}
\newcommand{\ba}{\mathbf{a}}
\newcommand{\bb}{\mathbf{b}}
\newcommand{\bp}{\mathbf{p}}
\newcommand{\bu}{\mathbf{u}}
\newcommand{\bv}{\mathbf{v}}
\newcommand{\bw}{\mathbf{w}}
\newcommand{\bz}{\mathbf{z}}
\newcommand{\slerp}{\operatorname{SLERP}}
\title[Analytic First Derivatives of SIDER Interpolation]{Analytic First Derivatives of SIDER Interpolation}
\author{Shingyu Leung}
\date{Source: arXiv:2606.13998v1 (CC BY 4.0)}
\begin{document}
\maketitle

\noindent\emph{Source and licence.} Analytic First Derivatives of SIDER Interpolation. Version arXiv:2606.13998v1 (CC BY 4.0), distributed under the Creative Commons Attribution 4.0 International licence, as recorded in the arXiv licence metadata for this exact version and shown on the abstract page https://arxiv.org/abs/2606.13998v1. Full rights notice: licenses/arxiv-2606.13998-CC-BY-4.0.txt.\par\smallskip

\noindent\textit{Editorial scope.} Counted unit: the Prop whose source label is \texttt{prop:sider-formula-tangent} in the article (source0.tex lines 820--837, source label \texttt{prop:sider-formula-tangent}), delivered together with the article's own complete original proof of it (source0.tex lines 839--908, one \texttt{proof} environment of 70 lines). No mathematics is written, repaired or re-derived: statement and proof are the article's own text line for line. Required context retained uncounted: eq:general-recursion (lines 417--428); lem:slerp-formula-tangent (lines 711--752). Every definition, display and label of those lines is retained unchanged. Omitted from this package and not redistributed: 1--416, 429--710, 753--819, 838--838, 909--1637. Nothing inside a retained excerpt is omitted or altered: every retained body is delivered exactly as the article's lines are written, so no omission receipt of any kind accompanies this package. Citations: the retained excerpts carry no citation command at all, so no bibliography entry of the article is transcribed and no reference is redistributed. Reference adaptation: none was needed and none was made; every reference inside the delivered unit resolves inside it, so no unresolved reference or citation is produced. Printed slips kept literal: no printed word, symbol, hypothesis or number of the counted statement or of its proof was altered. Layout and class: the article's own class is \texttt{article} with packages including grffile, latexsym, amssymb, enumerate, amsmath, epsfig, amsthm, geometry, setspace, color, graphics, graphicx, algorithm2e, empheq; the wrapper is the lane's pinned \texttt{amsart} editorial wrapper at 10pt on A4, which restates the theorem-like environments the retained text needs and adds the article's own macro definitions that the retained text needs (\texttt{\textbackslash bA}, \texttt{\textbackslash bB}, \texttt{\textbackslash bD}, \texttt{\textbackslash ba}, \texttt{\textbackslash bb}, \texttt{\textbackslash bp}, \texttt{\textbackslash bu}, \texttt{\textbackslash bv}, \texttt{\textbackslash bw}, \texttt{\textbackslash bz}, \texttt{\textbackslash slerp}), with extra wrapper packages (bm, bbm, dsfont, xspace, cancel, mathtools). The delivered numbering is the unit's own. Assets: no retained span contains a figure environment, a graphics-inclusion command, a PDF-page inclusion, a table or a TikZ picture; no image file is copied into this package, the package declares no asset dependency and no image file is redistributed with it. Licence: version arXiv:2606.13998v1 of this article is distributed under the Creative Commons Attribution 4.0 International licence, as recorded in the arXiv licence metadata for this exact version and shown on the abstract page https://arxiv.org/abs/2606.13998v1. A full rights notice is delivered at licenses/arxiv-2606.13998-CC-BY-4.0.txt. Characters: every retained excerpt is pure ASCII, so the delivered engine, XeLaTeX with its default Latin Modern text font, reports no missing character.
\begin{equation}
\label{eq:general-recursion}
\boxed{
        \bD_n[\bp_0,\ldots,\bp_n](t)
        =
        \mathcal{D}\slerp
        \left(
        \bA(t),\bB(t),t;
        \bA'(t),\bB'(t),1
        \right).
}
\end{equation}

\begin{lemma}[Tangency preservation of the SLERP derivative formula]\label{lem:slerp-formula-tangent}
Let \(\ba,\bb\in\mathbb{S}^2\) be non-antipodal, and let \(\bu\in T_{\ba}\mathbb{S}^2\), \(\bv\in T_{\bb}\mathbb{S}^2\).  Let \(\mu\in\mathbb{R}\) be the rate of change of the interpolation parameter.  Define
\[
        \theta=\arccos(\ba\cdot\bb),\qquad
        \alpha=\frac{\sin((1-\lambda)\theta)}{\sin\theta},
        \qquad
        \beta=\frac{\sin(\lambda\theta)}{\sin\theta},
\]
and set
\[
        \bz=\slerp(\ba,\bb;\lambda)=\alpha\ba+\beta\bb .
\]
Let
\[
        \dot\theta
        =
        -\frac{\bu\cdot\bb+\ba\cdot\bv}{\sin\theta},
\]
and define
\[
        \dot\alpha=\alpha_\lambda \mu+\alpha_\theta\dot\theta,
        \qquad
        \dot\beta=\beta_\lambda \mu+\beta_\theta\dot\theta .
\]
The algebraic SLERP derivative formula
\[
        \bw
        =
        \mathcal{D}\slerp(\ba,\bb,\lambda;\bu,\bv,\mu)
        =
        \alpha\bu+\beta\bv+\dot\alpha\,\ba+\dot\beta\,\bb
\]
satisfies
\[
        \bz\cdot\bw=0 .
\]
Consequently,
\[
        \mathcal{D}\slerp(\ba,\bb,\lambda;\bu,\bv,\mu)
        \in T_{\slerp(\ba,\bb;\lambda)}\mathbb{S}^2 .
\]
\end{lemma}

\begin{prop}[Tangency preservation of the recursive SIDER derivative formula]\label{prop:sider-formula-tangent}
Assume that all SLERP operations appearing in the SIDER construction are nonsingular.  Let
\[
        \mathbf{S}_n(t)=\mathbf{S}_n[\bp_0,\ldots,\bp_n](t)
\]
be the SIDER-\(n\) reconstruction, and let
\[
        \bD_n(t)=\bD_n[\bp_0,\ldots,\bp_n](t)
\]
be the vector produced by the recursive derivative formula \eqref{eq:general-recursion}, with the SIDER2 derivative used as the base case.  Then, in exact arithmetic,
\[
        \mathbf{S}_n(t)\cdot\bD_n(t)=0
\]
for every admissible \(t\).  Equivalently,
\[
        \bD_n(t)\in T_{\mathbf{S}_n(t)}\mathbb{S}^2 .
\]
\end{prop}

\begin{proof}
The proof is by induction on the order \(n\), but the induction is applied to the derivative formula rather than to an already-known smooth sphere-valued curve.

For the base case \(n=2\), the SIDER2 construction is
\[
        \mathbf{S}_2(t)=\slerp(\mathbf{A}_2(t),\mathbf{B}_2(t);t),
\]
where \(\mathbf{A}_2(t)\) and \(\mathbf{B}_2(t)\) are themselves SLERP curves.  The derivatives \(\mathbf{A}_2'(t)\) and \(\mathbf{B}_2'(t)\) are obtained from the fixed-endpoint SLERP derivative formula.  By Lemma~\ref{lem:slerp-formula-tangent}, these derivative vectors satisfy
\[
        \mathbf{A}_2(t)\cdot\mathbf{A}_2'(t)=0,
        \qquad
        \mathbf{B}_2(t)\cdot\mathbf{B}_2'(t)=0 .
\]
Therefore the two branch derivative vectors supplied to the outer SLERP derivative formula are tangent to their respective branch values.  Applying Lemma~\ref{lem:slerp-formula-tangent} once more to the outer SLERP gives
\[
        \mathbf{S}_2(t)\cdot\bD_2(t)=0 .
\]
Thus the derivative vector produced by the SIDER2 formula lies in \(T_{\mathbf{S}_2(t)}\mathbb{S}^2\).

Assume now that the statement holds for SIDER-\((n-1)\).  For SIDER-\(n\), write
\[
        \mathbf{S}_n(t)=\slerp(\mathbf{A}(t),\mathbf{B}(t);t),
\]
where
\[
        \mathbf{A}(t)
        =
        \mathbf{S}_{n-1}[\bp_0,\ldots,\bp_{n-1}](g_n(t)),
        \qquad
        \mathbf{B}(t)
        =
        \mathbf{S}_{n-1}[\bp_1,\ldots,\bp_n](h_n(t)).
\]
The recursive derivative formula uses
\[
        \mathbf{A}'_{\rm alg}(t)
        =
        \frac{n}{n-1}
        \bD_{n-1}[\bp_0,\ldots,\bp_{n-1}](g_n(t)),
\]
and
\[
        \mathbf{B}'_{\rm alg}(t)
        =
        \frac{n}{n-1}
        \bD_{n-1}[\bp_1,\ldots,\bp_n](h_n(t)).
\]
By the induction hypothesis,
\[
        \mathbf{A}(t)\cdot \mathbf{A}'_{\rm alg}(t)=0,
        \qquad
        \mathbf{B}(t)\cdot \mathbf{B}'_{\rm alg}(t)=0,
\]
because multiplication by the scalar factor \(n/(n-1)\) does not change tangency.  Thus the inputs passed to the outer SLERP derivative formula satisfy precisely the hypotheses of Lemma~\ref{lem:slerp-formula-tangent}.  Therefore
\[
        \mathbf{S}_n(t)\cdot
        \mathcal{D}\slerp
        \left(
        \mathbf{A}(t),\mathbf{B}(t),t;
        \mathbf{A}'_{\rm alg}(t),\mathbf{B}'_{\rm alg}(t),1
        \right)
        =
        0 .
\]
By definition, the vector on the second factor is exactly the recursively computed derivative \(\bD_n(t)\).  Hence
\[
        \mathbf{S}_n(t)\cdot\bD_n(t)=0 .
\]
This completes the induction.
\end{proof}

\end{document}
